\documentclass[10pt,a4paper]{amsart}
\usepackage[margin=19mm]{geometry}
\usepackage{amsmath,amssymb,mathtools}
\usepackage[scr=rsfs,cal=euler]{mathalfa}
\usepackage{xfrac}
\usepackage[unicode,hidelinks,bookmarks=false]{hyperref}
\newcommand{\ie}{i.e.\ }
\newcommand{\cf}{cf.\ }
\newcommand{\resp}{respectively\ }
\newcommand{\Subsection}{\subsection}
\providecommand{\nobreakdash}{\nobreak\hbox{-}}
\setlength{\emergencystretch}{2em}
\multlinegap0pt
\usepackage{amssymb}
\usepackage{mathtools}
\usepackage[shortlabels]{enumitem}
\newtheorem{lemma}{Lemma}
\newtheorem{proposition}{Proposition}
\newcommand{\R}{\mathbb{R}}
\newcommand{\Z}{\mathbb{Z}}
\newcommand{\e}{\mathrm{e}}
\title{Admissibility of H\"ormander--Bernhardsson extremal zeros}
\author{Khai-Hoan Nguyen-Dang}
\date{Source: arXiv:2602.10497v1 (CC BY 4.0)}
\begin{document}
\maketitle

\noindent\emph{Source and licence.} Admissibility of H\"ormander--Bernhardsson extremal zeros. Version arXiv:2602.10497v1 (CC BY 4.0), distributed by arXiv under the Creative Commons Attribution 4.0 International licence, http://creativecommons.org/licenses/by/4.0/, as recorded in the arXiv licence metadata for this exact version and shown on the abstract page https://arxiv.org/abs/2602.10497v1. Full licence notice: \texttt{licenses/arxiv-2602.10497-CC-BY-4.0.txt}.\par\smallskip

\noindent\textit{Editorial scope.} Counted unit: the article's proposition prop:Fr-asymp (source0.tex lines 956--966). The article's complete original proof of it is delivered with it (source0.tex lines 968--1057, one \texttt{proof} environment of 90 lines). No mathematics is written, repaired or re-derived: the statement and the proof are the article's own lines, in the article's own notation, and no retained line is substituted or re-flowed.  Retained uncounted as the source-local context the proof needs: lines 724--726; lines 766--773; lines 785--791; lines 851--871; lines 857--863.  Nothing was omitted from any retained span, so the package carries no omission record.  No retained line contains a citation, so the wrapper carries no bibliography and no bibliography entry.  No retained line contains a figure environment, an image-inclusion command or drawing code, so no image is delivered and the package declares no asset dependency.  Editorial formatting only, in addition: an AMS \texttt{amsart} preamble that restates the article's own declarations and macros used by the retained lines, namely the environments \texttt{lemma}, \texttt{proposition} on separate counters, together with the standard packages those lines need.  The retained environments print in this document as Lemma 1, Proposition 1 in order of appearance, whereas the article prints the same items with the numbering of its own full document.  The complete native source of the article as distributed in the arXiv e-print for this exact version is bundled unchanged as \texttt{sources/194403/source0.tex}.
\begin{equation}\label{eq:def-Fr}
F_r(t):=\sum_{n\ge 1}q(u_n)^r\,\e^{-a_n^2 t}.
\end{equation}

\begin{lemma}[Geometric bound for the coefficients $c_{r,m}$]\label{lem:crm-bound}
Fix $r\ge 0$ and choose any $\rho$ with $r_0<\rho<4$. Then there exists $C_{r,\rho}>0$
such that
\begin{equation}\label{eq:crm-bound}
|c_{r,m}|\le C_{r,\rho}\,\rho^{-m}\qquad(m\ge 0).
\end{equation}
In particular, the series $\sum_{m\ge 0}|c_{r,m}|\,r_0^m$ converges.
\end{lemma}

\begin{proposition}[Expansion of $F_r(t)$ in terms of $\Theta_m(t)$]\label{prop:Fr-Theta}
For every $r\ge 0$ and $t>0$,
\begin{equation}\label{eq:Fr-Theta}
F_r(t)=\sum_{m=0}^\infty c_{r,m}\,\Theta_m(t),
\end{equation}
and the series on the right converges absolutely.
\end{proposition}

\begin{proposition}[Full asymptotic expansion of $\Theta_m(t)$]\label{prop:Theta-asymp}
For each integer $m\ge 0$, the function $\Theta_m(t)$ admits a full asymptotic expansion
as $t\to 0^+$ with exponents in $\tfrac12\Z$.

More precisely, fix an integer $K\ge 0$ and set $s_K:=-K-\tfrac12$. Then there exist
constants $A_{m,j}$ ($0\le j\le K$) such that, as $t\to 0^+$,
\begin{equation}\label{eq:Theta-K}
\Theta_m(t)
=\frac{\sqrt{\pi}}{2}\,\mathbf{1}_{\{m=0\}}\,t^{-1/2}
\;+\;\sum_{j=0}^{K}A_{m,j}\,t^{j}
\;+\;\mathbf{1}_{\{1\le m\le K\}}\frac12\,\Gamma\!\Bigl(\frac12-m\Bigr)\,t^{m-1/2}
\;+\;O\!\bigl(t^{K+1/2}\bigr).
\end{equation}
In particular,
\[
\Theta_0(t)=O(t^{-1/2}),\qquad
\Theta_m(t)=\Theta_m(0)+O(t^{1/2})\ (m=1),\qquad
\Theta_m(t)=\Theta_m(0)+O(t)\ (m\ge 2),
\]
and hence $\Theta_m(t)=O(1)$ for every $m\ge 1$.
\end{proposition}

\begin{equation}
\Theta_m(t)
=\frac{\sqrt{\pi}}{2}\,\mathbf{1}_{\{m=0\}}\,t^{-1/2}
\;+\;\sum_{j=0}^{K}A_{m,j}\,t^{j}
\;+\;\mathbf{1}_{\{1\le m\le K\}}\frac12\,\Gamma\!\Bigl(\frac12-m\Bigr)\,t^{m-1/2}
\;+\;O\!\bigl(t^{K+1/2}\bigr).
\end{equation}

\begin{proposition}[Full asymptotic expansion of $F_r(t)$]\label{prop:Fr-asymp}
For each integer $r\ge 0$, the function $F_r(t)$ defined in \eqref{eq:def-Fr}
admits a full asymptotic expansion as $t\to 0^+$ with exponents in $\tfrac12\Z$.

More precisely, for every integer $K\ge 0$ there exist constants $B_{r,j}$,
$j\in\{-\tfrac12,0,\tfrac12,1,\dots,K\}$, such that
\begin{equation}\label{eq:Fr-K}
F_r(t)=\sum_{j\in\{-\tfrac12,0,\tfrac12,1,\dots,K\}}B_{r,j}\,t^{j}+O(t^{K+1/2})
\qquad(t\to 0^+).
\end{equation}
\end{proposition}

\begin{proof}
Fix $r\ge 0$ and $K\ge 0$. By Proposition~\ref{prop:Fr-Theta},
\[
F_r(t)=\sum_{m=0}^\infty c_{r,m}\,\Theta_m(t),
\]
with absolute convergence for every $t>0$. For each $m\ge0$, apply
Proposition~\ref{prop:Theta-asymp} (with the same $K$) and write
\[
\Theta_m(t)=\mathcal{P}_{m,K}(t)+R_{m,K}(t),
\]
where $\mathcal{P}_{m,K}(t)$ is the truncated part from \eqref{eq:Theta-K} and
$R_{m,K}(t)=O(t^{K+1/2})$ as $t\to0^+$.

\medskip
\noindent\emph{Uniform control of the remainder and summation over $m$.}
Set $M:=K+2$. Split
\[
\sum_{m\ge0}c_{r,m}R_{m,K}(t)=\sum_{m\le M}c_{r,m}R_{m,K}(t)+\sum_{m>M}c_{r,m}R_{m,K}(t).
\]
The first sum is finite, hence $O(t^{K+1/2})$.

For the tail $m>M$, we use the remainder representation coming from the contour shift
in the proof of Proposition~\ref{prop:Theta-asymp}. Namely, if we write
$s_K:=-K-\tfrac12$, then (after subtracting residues up to order $K$)
\begin{equation}\label{eq:Rmk-integral}
R_{m,K}(t)=\frac{1}{2\pi i}\int_{\Re s=s_K}\Gamma(s)\,t^{-s}\,S_m(s)\,ds,
\end{equation}
where
\[
S_m(s)=\sum_{n\ge1}\bigl(n+\tfrac12\bigr)^{-2(m+s)}.
\]
Since $x\mapsto x^{-\alpha}$ is decreasing on $[a_1,\infty)$ for $\alpha>1$, we have the standard bound
\[
\sum_{n\ge1} a_n^{-\alpha}
\le a_1^{-\alpha}+\int_{a_1}^{\infty}x^{-\alpha}\,dx
= a_1^{-\alpha}+\frac{a_1^{-\alpha+1}}{\alpha-1}.
\]
Apply this with $\alpha:=2(m+\Re s)$. On the line $\Re s=s_K=-K-\tfrac12$ and for $m\ge K+2$,
we have $\alpha=2m-2K-1\ge3$, hence $(\alpha-1)^{-1}\le \tfrac12$. Therefore
\[
|S_m(s)|
\le \sum_{n\ge1}a_n^{-2(m+\Re s)}
\ll_K a_1^{-2(m+\Re s)}+a_1^{-2(m+\Re s)+1}
\ll_K a_1^{-2m}.
\]
Since $a_1=\tfrac32$, this gives the geometric bound $|S_m(s)|\ll_K (4/9)^m$ uniformly in $\Im s$. Since $a_1^{-2}=(\tfrac23)^2=\tfrac49=:r_0$, this yields the uniform geometric bound
\begin{equation}\label{eq:Sm-geometric}
|S_m(s)|\ll_{K} r_0^{\,m}\qquad(\Re s=s_K,\ m>M),
\end{equation}
uniformly in $\Im s$. Inserting \eqref{eq:Sm-geometric} into \eqref{eq:Rmk-integral} gives, for
$t\in(0,1]$ and $m>M$,
\[
|R_{m,K}(t)|
\le \frac{t^{-s_K}}{2\pi}\Bigl(\sup_{\Re s=s_K}|S_m(s)|\Bigr)\int_{-\infty}^{\infty}|\Gamma(s_K+iy)|\,dy
\ll_{K} r_0^{\,m}\,t^{K+1/2},
\]
since $\int_{\R}|\Gamma(s_K+iy)|\,dy<\infty$ for fixed $s_K$.
Therefore, using Lemma~\ref{lem:crm-bound} (i.e.\ $\sum_{m\ge0}|c_{r,m}|r_0^{\,m}<\infty$),
\[
\sum_{m>M}|c_{r,m}||R_{m,K}(t)| \ll_K t^{K+1/2}\sum_{m>M}|c_{r,m}|r_0^{\,m}
=O(t^{K+1/2}).
\]
Hence
\[
\sum_{m\ge0}c_{r,m}R_{m,K}(t)=O(t^{K+1/2})\qquad(t\to0^+).
\]

\medskip
\noindent\emph{Summation of the truncated pieces.}
The function $\mathcal{P}_{m,K}(t)$ is a finite linear combination of the monomials
$t^{-1/2}$, $t^0,t^1,\dots,t^K$, and (when $1\le m\le K$) the half-integer monomial $t^{m-1/2}$.
Thus $\sum_{m\ge0}c_{r,m}\mathcal{P}_{m,K}(t)$ is a finite linear combination of the same exponents,
once we justify absolute convergence for the integer-power coefficients.

For the integer powers $t^j$ ($0\le j\le K$), Proposition~\ref{prop:Theta-asymp} gives
\[
A_{m,j}=\frac{(-1)^j}{j!}\,S_m(-j).
\]
If $m\ge j+2$, then $S_m(-j)=\sum_{n\ge1}a_n^{-2(m-j)}$ converges absolutely and the same estimate as
\eqref{eq:Sm-geometric} yields
\[
|S_m(-j)|\ll_j r_0^{\,m}.
\]
Hence $\sum_{m\ge0}|c_{r,m}A_{m,j}|<\infty$ for each $j$, because only finitely many $m<j+2$ remain and
$\sum_{m\ge0}|c_{r,m}|r_0^{\,m}<\infty$.

Collecting coefficients of the finitely many exponents
$j\in\{-\tfrac12,0,\tfrac12,1,\dots,K\}$ yields constants $B_{r,j}$ such that
\eqref{eq:Fr-K} holds.
\end{proof}

\end{document}
